\section{Results}

\subsection{Cross-Sectional Polarization Trend}

Table~\ref{tab:nominate-trend} reproduces the Voteview trend for context.

\begin{table}[h]
\centering\small
\caption{Mean $|\text{DW-NOMINATE dim1}|$ by decade and party.
Source: Voteview.com, downloaded 2026-05-09.}
\label{tab:nominate-trend}
\begin{tabular}{lcccc}
\toprule
Decade & Democrat & Republican & Combined & $N$ (member-congress) \\
\midrule
2001--2010 & 0.371 & 0.428 & 0.398 & 2{,}220 \\
2011--2020 & 0.386 & 0.483 & 0.432 & 2{,}229 \\
2021--2024 & 0.382 & 0.505 & 0.443 & 906 \\
\bottomrule
\end{tabular}
\end{table}

Republicans have polarized more rapidly (+0.077 abs.\ dim1 from 2000s to 2020s)
than Democrats (+0.011), consistent with the asymmetric polarization literature
\citep{hacker2006off}.

\subsection{Compactness and Ideological Extremism}

Table~\ref{tab:nominate-compactness} reports the OLS results.
The dependent variable is $|\text{DW-NOMINATE dim1}|$ for each member-congress.

\begin{table}[h]
\centering\small
\caption{OLS: Member ideological extremism ($|\text{DW-NOMINATE}|$) on district PP.
Standard errors clustered by member in parentheses.
$N = 5{,}355$ member-congress records.}
\label{tab:nominate-compactness}
\begin{tabular}{lcccc}
\toprule
 & (1) & (2) & (3) & (4) Full \\
\midrule
PP (per 0.10) & $-$0.091$^{**}$ & $-$0.079$^{*}$ & $-$0.074$^{*}$ & $-$0.071$^{*}$ \\
 & (0.031) & (0.029) & (0.028) & (0.027) \\
D pres.\ vote share & & & $-$0.312$^{***}$ & $-$0.298$^{***}$ \\
 & & & (0.042) & (0.041) \\
Republican & 0.063$^{***}$ & 0.058$^{***}$ & 0.058$^{***}$ & 0.051$^{***}$ \\
 & (0.009) & (0.009) & (0.009) & (0.009) \\
Member FE & No & Yes & Yes & Yes \\
State FE & No & No & No & Yes \\
Congress FE & No & No & No & Yes \\
$R^2$ & 0.16 & 0.52 & 0.59 & 0.68 \\
\bottomrule
\end{tabular}
\end{table}

\paragraph{Interpretation.}
A 0.10 PP improvement in district compactness predicts a 0.071 reduction in
$|\text{DW-NOMINATE}|$ (column 4, full specification).
This effect is conditional on member fixed effects (the ``within-member''
effect of serving a more vs.\ less compact district) and district partisan
composition.

The effect size is modest relative to the party gap (+0.051 for Republicans
vs.\ Democrats) but economically meaningful: the PP difference between
\textsc{bisect} (0.361) and enacted gerrymanders (0.286) implies:
\[
  \Delta|\text{DW-NOMINATE}| = 0.071 \times (0.361 - 0.286) / 0.10 \approx -0.053.
\]

\subsection{Counterfactual Projection}

Under the \textsc{bisect} plan for the five study states (NC, WI, TX, FL, PA),
the projected mean $|\text{DW-NOMINATE}|$ is${}^\dagger$:

\begin{table}[h]
\centering\small
\caption{Projected DW-NOMINATE extremism under bisect plans.
${}^\dagger$Projection from within-member PP slope; not a causal estimate.}
\label{tab:nominate-projection}
\begin{tabular}{lcccc}
\toprule
State & Enacted PP (mean) & Bisect PP (mean) & $\Delta$PP & Projected $\Delta|\text{NOMINATE}|$ \\
\midrule
NC & 0.241 & 0.368 & $+$0.127 & $-$0.090 \\
WI & 0.266 & 0.394 & $+$0.128 & $-$0.091 \\
TX & 0.278 & 0.358 & $+$0.080 & $-$0.057 \\
FL & 0.312 & 0.371 & $+$0.059 & $-$0.042 \\
PA & 0.231 & 0.349 & $+$0.118 & $-$0.084 \\
\midrule
Mean & 0.266 & 0.368 & $+$0.102 & $-$0.072 \\
\bottomrule
\end{tabular}
\end{table}

The projected reduction in ideological extremism is $-0.072$ absolute
DW-NOMINATE points on average across the five states --- representing
approximately 13\% reduction from the current mean of $0.432$.
This projection is labeled preliminary; causal identification would require
an instrumental variable strategy (e.g., using redistricting cycle boundaries
as instruments for compactness change).
